% Cote 119: offprints, reports, courses, CV and letters (typed and printed part): archive cote 119, site page 227 (PDF page 228).
% typescript. Draft compared: none.
% AI-transcribed by Claude Opus 5.5 (claude-opus-5-5, adaptive thinking, effort high), Message Batches API, batch r119_full (msgbatch_01DwZNm5Y7YsKBWyAN4sFs8w), 2026-11-13; not yet checked by eye.
\scanpage{119/227}
\mspage{. 14 .}

sur ces mêmes schémas sur des bases générales, et dans
une certaine mesure ceux de D. Mumford [$*$] et de
S. Seshadri sur le passage au quotient, pour n'en
citer que quelques-uns.

d) \uline{Théories cohomologiques} :

1°) \uline{Cohomologie ``cohérente''} : résultats de finitude,
de comparaison avec la cohomologie formelle, cf.
[15, Chap. III]. Théorèmes de dualité et des résidus :
un exposé systématique de mes idées et résultats est
développé dans le séminaire de R. Hartshorne [$*$],
cf. aussi [18 bis].

2°) \uline{Cohomologie $\ell$-adique} : définition de la cohomo-%
logie étale, théorèmes de comparaison, de finitude,
de dimension cohomologique, de Lefschetz faible,
[16 SGA 4] ; théorèmes de dualité, formules de Lefschetz
et d'Euler-Poincaré, application aux fonctions L ,
[16, SGA 15].

3°) \uline{Cohomologie de De Rham} :[16 bis],[17 bis].

4°) \uline{Cohomologie cristalline} : quelques idées de
départ sont esquissées dans [18, Exp. IX], puis reprises
et systématisées dans la thèse de P. Berthelot [$*$],
et dans le travail de P. Berthelot et L. Illusie sur
les classes de Chern cristallines [$*$] .

e) \uline{Théorie du groupe fondamental} ([16, SGA 1, SGA 2,
SGA 7, Exp. I et II], [15 bis, n° 182], [19 bis]) :

D'un point de vue algébrico-géométrique,
tout était à faire, depuis la définition du groupe
fondamental d'une variété quelconque, en passant par
des propriétés ``de descente'' incluant des résultats
assez formels du type de van Kampen, jusqu'au calcul

\margin{227}
